%% kapa: exam=1 round=35 year=2024 period=2 form=A subject=appraisal_law subject_ko=감정평가관계법규 seq=22 number=22
%% source: Q-Net 감정평가사 기출문제 (cst00309) artlSeq=5213229 file=raw/1008/5213229/2교시 문제지 원본.pdf
%% conversion: kordoc 4.24.51 -> tools/gen1.py (auto); review=proofread-sonnet; flags=none
\begin{question}{22}
국유재산법령상 행정재산의 사용허가에 관한 설명으로 옳은 것은?
\begin{choices}
  \item 사용허가를 받은 자는 허가기간이 끝난 경우에는 중앙관서의 장이 미리 상태의 변경을 승인하였더라도 그 재산을 원래 상태대로 반환하여야 한다.
  \item 경작용으로 실경작자에게 사용허가를 하는 경우에는 수의의 방법으로 사용허가를 받을 자를 결정할 수 없다.
  \item 중앙관서의 장은 사용허가를 받은 자가 해당 재산의 보존을 게을리한 경우 그 허가를 철회할 수 있다.
  \item 사용허가에 관하여는 「국유재산법」에서 정한 것을 제외하고는 「민법」의 규정을 준용한다.
  \item 사용허가를 받은 자가 그 재산에 대하여 유지ㆍ보수 외의 시설을 설치하려는 때에는 총괄청의 허가를 받아야 한다.
\end{choices}
\end{question}
\answer{3}
